\documentclass[a4paper,11pt]{article}
\usepackage{exam-style-341}
\examinehead{341 农业综合知识三 · 数据库分章练习}{第01章\quad 基本概念与数据模型}
\begin{document}

\unitcover{341 农业综合知识三 · 数据库分章练习}{第01章\quad 基本概念与数据模型}{练习（共 100 分）}{约 60 分钟}{数据、数据库、数据库系统与 DBMS；数据管理三个发展阶段；数据库三级模式与两级映像；数据逻辑独立性与物理独立性；数据模型三要素；概念模型与 E-R 图（实体、属性、联系、1:1 / 1:n / m:n）；过程语言与非过程语言；DBMS 的主要功能。}

\begin{knowmap}
\textbf{概念链：}数据（记载事物的物理符号）$\rightarrow$ 数据库（长期存储、有组织、可共享的数据集合）$\rightarrow$ DBMS（管理数据的系统软件）$\rightarrow$ 数据库系统（数据库 + DBMS + 应用程序 + DBA + 硬件）\\[2pt]
\textbf{数据管理三阶段：}人工管理 $\rightarrow$ 文件系统 $\rightarrow$ 数据库系统；数据独立性由无到有、由弱到强，冗余度由高到低\\[2pt]
\textbf{三级模式：}外模式（用户视图，多个）$\rightarrow$ 模式（全局逻辑结构，唯一）$\rightarrow$ 内模式（物理存储结构，唯一）\\[2pt]
\textbf{两级映像：}外模式/模式映像 $\Rightarrow$ 逻辑独立性；模式/内模式映像 $\Rightarrow$ 物理独立性\\[2pt]
\textbf{数据模型三要素：}数据结构（静态特性）· 数据操作（动态特性）· 完整性约束；概念模型用 E-R 图表示，三要素为实体、属性、联系\\[2pt]
\textbf{语言特性：}过程语言（指出“怎么做”，如 C）与非过程语言（只说明“做什么”，如 SQL）
\end{knowmap}

\begin{keyterms}
\textbf{数据}：描述事物的符号记录，是数据库中存储的基本对象，有数值、文字、图形、声音等形式。\\
\textbf{数据库}：长期存储在计算机内、有组织、可共享的大量数据的集合，按一定的数据模型组织、描述和存储。\\
\textbf{数据库管理系统}：位于用户与操作系统之间、用于科学地组织和存储数据、高效地获取和维护数据的大型系统软件。\\
\textbf{数据库系统}：由数据库、数据库管理系统、应用程序、数据库管理员和硬件等构成的整体，其核心是 DBMS。\\
\textbf{外模式}：又称子模式或用户模式，是数据库用户能够看见和使用的局部数据的逻辑结构。\\
\textbf{模式}：又称逻辑模式，是数据库中全体数据的全局逻辑结构和特征的描述，一个数据库只有一个模式。\\
\textbf{内模式}：又称存储模式，是数据物理结构和存储方式的描述，一个数据库只有一个内模式。\\
\textbf{外模式/模式映像}：定义外模式与模式之间的对应关系；当模式改变时修改该映像，可使外模式及应用程序保持不变。\\
\textbf{模式/内模式映像}：定义模式与内模式之间的对应关系；当内模式的存储结构改变时修改该映像，可使模式保持不变。\\
\textbf{数据逻辑独立性}：当模式改变时，通过修改外模式/模式映像使外模式不变，从而应用程序不必修改。\\
\textbf{数据物理独立性}：当内模式改变时，通过修改模式/内模式映像使模式不变，从而应用程序不必修改。\\
\textbf{数据模型}：对现实世界数据特征的抽象，由数据结构、数据操作和完整性约束三要素组成。\\
\textbf{概念模型}：按用户的观点对数据和信息建模的模型，是现实世界到机器世界的中间层次，常用 E-R 图表示。\\
\textbf{非过程语言}：只说明“做什么”而不必指出“怎么做”的语言，SQL 是典型的非过程化语言。
\end{keyterms}

\begin{contrast}{数据 / 数据库 / 数据库管理系统 / 数据库系统（每行依次为：定义·组成或对象·边界与常见误解）}
\citem{数据}{描述事物的符号记录}{数值、文字、图形、声音、图像等}{数据与其语义密不可分，单看符号不能确定含义}
\citem{数据库}{长期存储、有组织、可共享的大量数据集合}{按数据模型组织、描述和存储的数据}{数据库不等于 DBMS，也不等于数据库系统}
\citem{数据库管理系统}{位于用户与操作系统之间的数据管理软件}{数据定义、数据操纵、运行管理、建立与维护}{是数据库系统的核心，但不是数据库本身}
\citem{数据库系统}{数据库、DBMS、应用程序、DBA、硬件的整体}{人、机、数据与软件的集合}{范围最大，包含数据库和 DBMS}
\end{contrast}

\begin{contrast}{外模式 / 模式 / 内模式（每行依次为：面向对象·数量·作用）}
\citem{外模式}{数据库用户（应用程序）}{一个数据库可有多个外模式}{描述用户可见的局部数据逻辑结构，是用户视图}
\citem{模式}{全体数据}{一个数据库只有一个模式}{描述全局逻辑结构，是三级模式的核心与中间层}
\citem{内模式}{物理存储}{一个数据库只有一个内模式}{描述数据的物理结构、存储方式与存取路径}
\end{contrast}

\begin{contrast}{逻辑独立性 / 物理独立性（每行依次为：含义·由哪级映像保证·模式变化时的影响）}
\citem{逻辑独立性}{模式改变而外模式与应用程序不变}{外模式/模式映像}{修改模式后只需改映像，应用程序不必修改}
\citem{物理独立性}{内模式改变而模式与应用程序不变}{模式/内模式映像}{改变存储结构或存取方法后只需改映像}
\end{contrast}

\begin{contrast}{过程语言 / 非过程语言（每行依次为：含义·代表语言·特点）}
\citem{过程语言}{既要说明做什么，也要指出怎么做}{C、Pascal 等}{逐条描述处理步骤，与数据的存储结构关系密切}
\citem{非过程语言}{只说明做什么，不必指出怎么做}{SQL}{面向集合、使用方便、可移植性好，是关系数据库的标准语言}
\end{contrast}

\qtypename{一、单项选择题}{10 题，每题 2 分，共 20 分}
\q 数据库系统的核心是（\quad）
\choicesline{A. 数据库}{B. 数据库管理系统}{C. 应用程序}{D. 数据库管理员}

\q 在数据库三级模式结构中，描述数据库中全体数据的全局逻辑结构的是（\quad）
\choicesline{A. 外模式}{B. 模式}{C. 内模式}{D. 存储模式}

\q 当数据库的模式发生改变时，通过修改（\quad）可以使外模式保持不变，从而保证数据的逻辑独立性。
\choicesline{A. 外模式/模式映像}{B. 模式/内模式映像}{C. 外模式}{D. 内模式}

\q 关于三级模式的数量，下列说法正确的是（\quad）
\choicesline{A. 模式可以有多个}{B. 内模式可以有多个}{C. 外模式可以有多个}{D. 三者都只能有一个}

\q 数据模型的三要素是（\quad）
\choices{数据结构、数据操作、完整性约束}{实体、属性、联系}{外模式、模式、内模式}{数据库、DBMS、数据库系统}

\q 用 E-R 图描述现实世界中的事物及其联系，这种模型属于（\quad）
\choicesline{A. 概念模型}{B. 逻辑模型}{C. 物理模型}{D. 关系模型}

\q 在数据管理的三个发展阶段中，数据独立性最高、冗余度最低的是（\quad）
\choicesline{A. 人工管理阶段}{B. 文件系统阶段}{C. 数据库系统阶段}{D. 三个阶段相同}

\q SQL 语言属于（\quad）
\choicesline{A. 过程语言}{B. 非过程语言}{C. 机器语言}{D. 汇编语言}

\q 下列不属于数据库管理系统主要功能的是（\quad）
\choicesline{A. 数据定义}{B. 数据操纵}{C. 数据库的运行管理}{D. 编译并运行 C 源程序}

\q 现实世界中的事物经过人的认识形成信息世界，再由信息世界转换为机器世界。E-R 图描述的是（\quad）
\choicesline{A. 现实世界}{B. 信息世界}{C. 机器世界}{D. 存储世界}

\qtypename{二、多项选择题}{5 题，每题 4 分，共 20 分；每题有两个或两个以上正确项}
\q 数据库系统通常由（\quad）等部分组成。
\choicesline{A. 数据库}{B. 数据库管理系统}{C. 应用程序}{D. 数据库管理员}

\q 与文件系统相比，数据库系统的主要特点有（\quad）
\choices{数据结构化且由 DBMS 统一管理}{数据共享性高、冗余度低}{数据独立性高}{数据完全由应用程序自行管理}

\q 关于数据库三级模式，下列说法正确的有（\quad）
\choices{一个数据库只有一个模式}{一个数据库可以有多个外模式}{一个数据库只有一个内模式}{模式是全体数据的全局逻辑结构}

\q 下列属于数据模型组成要素的有（\quad）
\choicesline{A. 数据结构}{B. 数据操作}{C. 完整性约束}{D. 数据独立性}

\q 关于 E-R 图，下列说法正确的有（\quad）
\choicesline{A. 实体用矩形表示}{B. 属性用椭圆表示}{C. 联系用菱形表示}{D. 联系只有 1:1 一种类型}

\qtypename{三、判断题}{5 题，每题 2 分，共 10 分；正确的画“√”，错误的画“×”并写出改正后的表述}
\q 数据库是长期存储在计算机内、有组织、可共享的大量数据的集合。\quad（\quad）\quad 改正：\fixblank

\q 数据库系统的核心是数据库。\quad（\quad）\quad 改正：\fixblank

\q 一个数据库可以有多个模式，也可以有多个内模式。\quad（\quad）\quad 改正：\fixblank

\q 模式/内模式映像保证了数据的逻辑独立性。\quad（\quad）\quad 改正：\fixblank

\q 在 E-R 图中，一个联系只能涉及两个实体。\quad（\quad）\quad 改正：\fixblank

\qtypename{四、填空题}{5 题，每题 2 分，共 10 分}
\q 数据管理技术的发展经历了 \fillblank{}、\fillblank{} 和数据库系统三个阶段。

\q 数据模型的三要素是数据结构、\fillblank{} 和 \fillblank{}。

\q 数据库的三级模式分别是外模式、\fillblank{} 和 \fillblank{}；两级映像分别是外模式/模式映像和 \fillblank{}。

\q E-R 图的三要素是实体、\fillblank{} 和联系；两个实体之间的联系类型有 1:1、1:n 和 \fillblank{} 三种。

\q 数据独立性包括数据的 \fillblank{} 独立性和 \fillblank{} 独立性。

\qtypename{五、简答题}{2 题，每题 8 分，共 16 分}
\q 简述数据库系统的三级模式结构及两级映像，并说明两级映像如何保证数据的逻辑独立性和物理独立性。\anslines{2.2cm}

\q 简述数据模型的三要素，并说明概念模型的作用以及 E-R 图的三要素。\anslines{2.2cm}

\qtypename{六、综合题}{2 题，每题 12 分，共 24 分}
\q 某农业合作社拟建立作物种植与订单管理数据库，经需求分析得到如下信息：
\begin{itemize}[leftmargin=1.6em, itemsep=1pt]
  \item 农户：农户编号、姓名、联系电话；
  \item 作物：作物编号、作物名称、单价（元/千克）；
  \item 订单：订单号、订单日期；
  \item 一个农户可以种植多种作物，一种作物也可以由多个农户种植，种植联系上需要记录面积（亩）；
  \item 一个农户可以下多个订单，一个订单只属于一个农户；
  \item 一个订单可以订购多种作物，一种作物也可以出现在多个订单中，订购联系上需要记录数量（千克）。
\end{itemize}

请完成：
\begin{enumerate}[label=（\arabic*）, leftmargin=2.4em, itemsep=1pt]
  \item 指出该应用中的实体、属性和联系，并写明每个联系的类型（1:1、1:n 或 m:n）。（4 分）
  \item 画出该应用的 E-R 图（实体、属性、联系及其类型）。（4 分）
  \item 将 E-R 图转换为关系模式，指出每个关系模式的主码和外码，并说明 1:n 联系与 m:n 联系的转换规则。（4 分）
\end{enumerate}
\anslines{4.5cm}

\q 某商业银行使用数据库系统管理业务数据，现发生以下三件事：
\begin{enumerate}[label=（\arabic*）, leftmargin=2.4em, itemsep=1pt]
  \item 因新增“客户风险等级”业务，DBA 在全局逻辑结构中增加了相应数据项；
  \item 因设备更新，DBA 把数据文件从机械硬盘迁移到固态硬盘，并改变了存储组织和存取路径；
  \item 为满足不同部门需要，DBA 分别为信贷部、柜面部分别建立了只包含本部门所需数据的视图。
\end{enumerate}

请分别分析：上述每件事分别改变了哪一级模式、是否需要修改映像、应修改哪一个映像、应用程序是否需要修改，并说明各体现了哪种数据独立性。最后归纳三级模式与两级映像的关系。（12 分）
\anslines{4.5cm}

\end{document}
